% BHU PYQ -- Manually transcribed & error-corrected
% Paper: BSC-04A: Elements of Geology - II
% Exam: B.Sc. (Hons.) Semester IV Examination, 2022-23

\documentclass[11pt,a4paper]{article}
\usepackage[a4paper,top=2.5cm,bottom=2.5cm,left=2.8cm,right=2.8cm,headheight=14pt]{geometry}
\usepackage{amsmath,amssymb}
\usepackage[T1]{fontenc}
\usepackage[utf8]{inputenc}
\usepackage{lmodern,microtype,fancyhdr,enumitem,hyperref,lastpage,parskip}

\hypersetup{colorlinks=true,urlcolor=black,linkcolor=black,
  pdftitle={BSC-04A Elements of Geology - II -- BHU B.Sc. Sem IV 2022-23}}
\pagestyle{fancy}\fancyhf{}
\renewcommand{\headrulewidth}{0.4pt}\renewcommand{\footrulewidth}{0.4pt}
\fancyhead[L]{\small\textit{Banaras Hindu University}}
\fancyhead[R]{\small\textit{BSC-04A: Elements of Geology - II}}
\fancyfoot[C]{\small Page \thepage\ of \pageref{LastPage}}

\newlist{parts}{enumerate}{2}
\setlist[parts,1]{label=(\alph*),leftmargin=*,itemsep=5pt,topsep=3pt}
\newlist{romanparts}{enumerate}{2}
\setlist[romanparts,1]{label=(\roman*),leftmargin=*,itemsep=5pt,topsep=3pt}
\newcommand{\pts}[1]{\hfill\textit{[#1]}}

\begin{document}
\begin{center}
  {\large\bfseries BANARAS HINDU UNIVERSITY}\\[2pt]
  {\normalsize B.Sc. (Hons.) --- Semester IV Examination, 2022-23}\\[6pt]
  \rule{0.8\linewidth}{0.5pt}\\[6pt]
  {\large\bfseries Paper: BSC-04A --- Elements of Geology - II}\\[4pt]
  {\normalsize Time: 3 Hours \hfill Maximum Marks: 70}\\[2pt]
  {\small Ref. No.: 058432}
\end{center}
\rule{\linewidth}{0.4pt}
\smallskip
\noindent\textit{Instructions: Answer \textbf{five} questions in all, selecting at least \textbf{two} questions from each of Sections A and B. Section-C is compulsory. All questions carry equal marks.}
\bigskip

\bigskip\noindent{\large\bfseries SECTION --- A}
\rule{\linewidth}{0.4pt}\smallskip

\noindent\textbf{Question 1.} What is the significance of the geological time scale? Explain the representative life forms of different geological periods. \pts{14}

\medskip\noindent\textbf{Question 2.} Answer the following:
\begin{parts}
  \item What are the different stratigraphic units? Describe the difference between lithostratigraphic and chronostratigraphic units. \pts{7}
  \item Write an essay on the principles of stratigraphy. How are these principles helpful in establishing a geological record of the strata? \pts{7}
\end{parts}

\medskip\noindent\textbf{Question 3.} Answer the following:
\begin{parts}
  \item Discuss the geological differences between the Andaman and Lakshadweep Islands of India. \pts{5}
  \item Discuss briefly the desertification of Thar and mention the characteristic topography of deserts. \pts{5}
  \item Write a short note on the mass extinction event related to dinosaurs. \pts{4}
\end{parts}

\medskip\noindent\textbf{Question 4.} Answer the following:
\begin{parts}
  \item Discuss the difference between the eastern and western coastal plains of India. \pts{4}
  \item Discuss ``Present is the key to the past''. \pts{4}
  \item Name the important peninsular plateaus of India. \pts{3}
  \item Write a short note on the ``Law of lateral continuity''. \pts{3}
\end{parts}

\bigskip\noindent{\large\bfseries SECTION --- B}
\rule{\linewidth}{0.4pt}\smallskip

\noindent\textbf{Question 5.} Discuss the conditions for fossilization of organisms and various kinds/types of fossils. \pts{14}

\medskip\noindent\textbf{Question 6.} Define paleontology. Describe the scope of paleontology in relation to other sub-disciplines of geology. \pts{14}

\medskip\noindent\textbf{Question 7.} Write explanatory notes on the following:
\begin{parts}
  \item Siwalik vertebrate fauna. \pts{6}
  \item Traces of biological activity. \pts{4}
  \item Living fossil. \pts{4}
\end{parts}

\bigskip\noindent{\large\bfseries SECTION --- C}
\rule{\linewidth}{0.4pt}\smallskip

\noindent\textbf{Question 8.} Define the following in a maximum of five sentences each:
\begin{romanparts}
  \item Acme Zone \pts{2}
  \item Lagoon \pts{2}
  \item Barren Island \pts{2}
  \item Carbonization \pts{2}
  \item Mould and Cast \pts{2}
  \item Pseudo fossil \pts{2}
  \item Facies fossil \pts{2}
\end{romanparts}

\vfill\rule{\linewidth}{0.4pt}
\begin{center}\small
\href{https://bhu-exam.vercel.app/}{Click here to visit our website}\\[4pt]
\href{https://me-aryan.vercel.app/}{Click here to visit our contributor}\\[4pt]
\href{https://quashan-me.vercel.app/}{Click here to visit our contributor}
\end{center}
\end{document}
